% =====================================================================
%  6. MENACES SUR LA VALIDITÉ
% =====================================================================
\section{Menaces sur la validité}\label{sec:menaces}

Nous suivons la classification de Wohlin et al. \cite{wohlin2012},
complétée par les recommandations de reporting des études de cas de
Runeson et H{\"o}st \cite{runeson2009}. Aucune de ces menaces n’est
résolue~; toutes sont déclarées.

\subsection{Validité de construction}\label{sec:construction}

Les métriques mobilisées — nombre de modules, de tests, d’avertissements,
taille du plus gros fichier — sont des \emph{indicateurs} de structure et
de défense, pas des mesures directes de qualité. Le compte de tests ne
mesure ni la couverture ni la force des tests~: \emph{aucune mesure de
couverture de code n’a été collectée} durant l’étude, et l’appréciation
qualitative des extensions (« parité des données embarquées »,
« comportements aux limites ») est celle de l’agent auteur. Le compte de
modules dépend de la discipline de cartographie elle-même — l’instrument
et l’objet mesuré sont partiellement confondus. Enfin, la « parité
fonctionnelle 14/14 » est une audit de familles de fonctionnalités guidé
par le fichier de lecture de la base, pas une équivalence formelle~; elle
ne couvre ni les comportements non documentés, ni les caractéristiques non
fonctionnelles au-delà des tests.

\subsection{Validité interne}\label{sec:interne}

Il n’existe \emph{pas de condition témoin}~: aucun recodage parallèle de la
même base sans la discipline TRIADE n’a été conduit, et l’attribution des
améliorations au protocole plutôt qu’à l’attention renouvelée portée au
code est un jugement, non une mesure. La structure de l’étude est
auto-référentielle~: l’agent qui a écrit les spécifications a aussi écrit
le code, les tests, l’audit et le présent article — un biais
d’auto-confirmation que seule une évaluation indépendante pourrait
lever. Les trois comparatifs avant/après ont été \emph{sélectionnés}
comme les plus illustratifs~; les vagues ont aussi produit des
conversions moins spectaculaires, non rapportées. Les mesures de
performance ont été prises sur une infrastructure partagée à 2~vCPU, dans
des conditions favorables au meilleur cas (meilleur de cinq essais).

\subsection{Validité externe}\label{sec:externe}

C’est la menace majeure~: l’étude est \emph{monocasale} (N~$=$~1). Un
projet, un domaine (audio temps réel), une pile technologique
(C\#/.NET~8/Avalonia), un agent (un unique modèle, dans un unique
environnement), un développeur, et — pour le protocole lui-même — ses
propres auteurs, puisque la démonstration a servi à valider la méthode
qui l’encadrait. Aucun claim de généralisation n’est avancé~: les
mécanismes rapportés sont \emph{plausibles} dans d’autres contextes, pas
\emph{établis}. La réplication sur d’autres projets, piles et agents,
conduite par des tiers, est la condition de toute montée en généralité
\cite{runeson2009}. De même, l’applicabilité à des équipes humaines
nombreuses — dont les incitations, contrairement à celles d’un agent,
comprennent la vitesse de livraison individuelle — est une question
ouverte que le protocole n’a pas encore affrontée.

\subsection{Validité de conclusion}\label{sec:conclusion-val}

Les affirmations de performance reposent sur des micro-benchmarks en
meilleur cas sur infrastructure partagée~; elles établissent l’\emph{absence
d’un problème de performance à corriger}, pas des performances
représentatives de déploiements réels. Le « zéro régression » est borné
par l’étendue de la suite de tests et de l’audit de familles~; il ne
préjuge pas des comportements non couverts — l’épisode du verrouillage de
journal sous Windows (invisible sous Linux/macOS) rappelle qu’une suite
verte est une \emph{preuve conditionnelle}, conditionnelle à
l’environnement d’exécution. L’épisode v3.0.0~$\to$~v3.0.2, présenté comme
validation du mécanisme d’échec honnête, est un événement unique~; sa
généralité n’est pas démontrée.

\subsection{L’écart conceptuel~: refactorisation source ou génération
binaire native~?}\label{sec:ecart}

Une menace spécifique mérite un traitement distinct, car elle touche au
rapport entre la démonstration et l’objet annoncé du protocole. La
spécification TRIADE v0.3 promet, à terme, un développement \emph{direct
en code machine natif scellé} — la couche B comme suite d’octets écrits
via une notation assemblée et vérifiée par round-trip, le compilateur
disparaissant du pipeline. L’étude de cas n’a rien démontré de tel~: la
couche B y est instanciée comme \emph{code source managé}, compilé par une
chaîne standard (compilation .NET, exécution avec compilation à la volée),
et les mécanismes spécifiques au natif — notation, round-trip, verrou
croisé cryptographique par empreintes d’octets, patch à chaud par
trampoline, adresses garanties par remplissage — n’ont pas été exercés.

L’étude valide donc la \emph{couche de gouvernance} du protocole — les
trois artefacts, les dix règles, la boucle, la vérification-puis-écriture,
la cristallisation, l’arbitrage lexicographique — appliquée à une
instanciation source de la couche B. Elle ne valide ni la forme de
l’artefact, ni les mécanismes natifs, ni l’argument de performance qui
les motive~; ces points restent la charge des jalons M0 à M7, non
encore franchis. Symétriquement, la motivation \emph{agent-first} du
protocole — l’économie de contexte par édition ciblée via la carte —
repose sur des budgets cibles (40 / 3\,000 / 6\,000 tokens) qui sont des
\emph{objectifs de conception}, pas des mesures~: aucune instrumentation
du coût réel en tokens n’a été conduite durant l’étude. Nous posons ces
limites d’autant plus clairement que la tentation symétrique existe~:
conclure de l’étude à la validité du programme natif serait aussi
infondé que conclure de son cadre à l’inutilité de la gouvernance. Ce
que l’étude établit tient en une phrase~: \emph{la discipline de
gouvernance est applicable de bout en bout, à coût documenté, à une
application réelle complète, et les défauts qu’elle a convertis en
contrats sont ceux-là mêmes qui survivaient depuis trois ans à la
maintenance classique}.
